%=====================================================================
\chapter{Research Questions, Hypotheses and Constructs}\label{ch:questions}
%=====================================================================
\locator{1}{Part~1 --- Foundations of Inquiry in Computing}

\begin{outcomes}
By the end of this chapter you will be able to:
\begin{tight}
  \item \textbf{Classify} a research question into one of six types and name
    the designs each type admits.
  \item \textbf{Apply} the FINER criteria to reject or repair a candidate
    question.
  \item \textbf{Decompose} a question into PICOC elements and use them to build
    a search string.
  \item \textbf{Operationalise} a construct into variables, measures and an
    instrument, and say where the definition could be contested.
  \item \textbf{Construct} an alignment table linking every research question
    to its data and its planned analysis.
  \item \textbf{Decide} whether your study should state hypotheses at all.
\end{tight}
\end{outcomes}

\begin{prereq}
\chref{contribution} for the claim--evidence--falsification triple, and
\chref{philosophy} for why the question type constrains the method. This
chapter turns both into a procedure.
\end{prereq}

\begin{keyterms}
research question $\bullet$ FINER $\bullet$ PICOC $\bullet$ construct $\bullet$
operationalisation $\bullet$ variable $\bullet$ independent and dependent
variable $\bullet$ moderator $\bullet$ mediator $\bullet$ confounder $\bullet$
null and alternative hypothesis $\bullet$ directional hypothesis $\bullet$
alignment table
\end{keyterms}

\section{The Alignment Chain}

Everything in a competent study is a chain, and a study fails at its weakest
link rather than at its most sophisticated one. A perfectly executed
mixed-effects model cannot rescue a measure that does not capture the construct
the question asked about.

\armfigh{
\begin{tikzpicture}[
  node distance=0pt,
  lk/.style={rounded corners=3pt, draw=primary!55, fill=primary!7,
             text width=2.15cm, align=center, font=\footnotesize,
             minimum height=1.05cm, inner sep=3pt},
  q/.style={font=\scriptsize\itshape, text=accent!80!black, align=center,
            text width=2.3cm},
  ar/.style={-{Stealth[length=2.2mm]}, draw=secondary, semithick}]

  \foreach \i/\name in {0/{\textbf{Topic}}, 1/{\textbf{Question}},
                        2/{\textbf{Construct}}, 3/{\textbf{Variable}},
                        4/{\textbf{Measure}}, 5/{\textbf{Analysis}}}{
    \node[lk] (n\i) at (\i*2.45,0) {\name};}

  \foreach \i [evaluate=\i as \j using int(\i+1)] in {0,...,4}{
    \draw[ar] (n\i) -- (n\j);}

  \node[q, below=5pt of n1] {what exactly\\is unknown?};
  \node[q, below=5pt of n2] {what abstract\\thing is it about?};
  \node[q, below=5pt of n3] {what observable\\stands for it?};
  \node[q, below=5pt of n4] {how is it\\recorded?};
  \node[q, below=5pt of n5] {does it answer\\the question?};

  \draw[ar, dashed, draw=accent, thick] (n5.north) .. controls +(0,1.1) and +(0,1.1) ..
       node[above, font=\scriptsize, text=accent]
       {check backwards before collecting any data} (n1.north);
\end{tikzpicture}}
{The alignment chain. Build it left to right, then verify it right to left:
given this analysis, on this measure, of this variable, standing for this
construct --- have I answered the question I asked? Most fatal design errors
are visible on that backward pass, and only on it.}{alignchain}

\section{Six Kinds of Question}

The type of question you ask determines the designs available to you. Asking a
causal question and then running a correlational study is not a small
mismatch --- it is the flaw that makes a study unpublishable.

\begin{center}
\small
\begin{longtable}{@{}L{2.5cm}L{5.2cm}L{6.5cm}@{}}
\toprule
\textbf{Type} & \textbf{Form} & \textbf{Designs that can answer it} \\
\midrule
\endfirsthead
\toprule
\textbf{Type} & \textbf{Form} & \textbf{Designs that can answer it} \\
\midrule
\endhead
\bottomrule
\endfoot
Exploratory
  & What is going on here? What forms does X take?
  & Qualitative study, grounded theory, exploratory case study
    (\chref{qualitative}, \chref{casestudy}) \\
Descriptive
  & How prevalent is X? What is the distribution of Y?
  & Survey with probability sampling; repository mining; systematic mapping
    (\chref{survey}, \chref{quasiexp}, \chref{slr}) \\
Relational
  & Is X associated with Y?
  & Correlational study, observational repository analysis
    (\chref{quasiexp}) \\
Causal
  & Does X \emph{cause} a change in Y?
  & Controlled experiment; quasi-experiment with an identification strategy
    (\chref{experiments}, \chref{quasiexp}) \\
Design
  & Can an artefact with properties P be built, and what does building it
    teach?
  & Design science research with a matched evaluation
    (\chref{designscience}) \\
Evaluative / comparative
  & Which of A and B performs better under conditions C?
  & Benchmarking; controlled comparison; classifier comparison across datasets
    (\chref{benchmarking}, \chref{mleval}) \\
\end{longtable}
\end{center}

\begin{paperbox}[Matching method to question]
Steve Easterbrook, Janice Singer, Margaret-Anne Storey and Daniela Damian,
\emph{Selecting Empirical Methods for Software Engineering Research}, in
\emph{Guide to Advanced Empirical Software Engineering}, Springer,
2008~\cite{easterbrook2008selecting}.

\medskip
\textbf{Why it matters.} It is the clearest treatment of the point this chapter
turns on: the method is chosen by the question and the philosophical stance,
not by preference or familiarity. It also connects method choice back to the
paradigms of \chref{philosophy} explicitly, which most methods texts do not.

\medskip
\textbf{What to extract.} Their mapping from question type to method, and their
discussion of what each method cannot tell you. Write the second one into your
threats section before you start.
\end{paperbox}

\section{FINER: Five Tests a Question Must Pass}

A question can be well-formed and still be a bad choice for \emph{you}, in
\emph{this} degree, with \emph{these} resources. The FINER
criteria~\cite{hulley2013designing} catch that early, when changing course is
still cheap.

\begin{center}
\small
\begin{tabular}{@{}L{2.2cm}L{5.0cm}L{6.7cm}@{}}
\toprule
 & \textbf{Test} & \textbf{What kills a question here} \\
\midrule
\textbf{F}easible
  & Subjects, time, expertise, money, access
  & 300 developers needed; you can recruit 20. A dataset that requires an
    agreement no one will sign. A study needing 18 months when you have 9 \\
\textbf{I}nteresting
  & To the field, not only to you
  & Nobody outside your supervisor's office would change any belief on hearing
    the answer \\
\textbf{N}ovel
  & Adds to what is known
  & Answered in 2019 and again in 2022. Only \chref{slr} can tell you this \\
\textbf{E}thical
  & Approvable, and defensible
  & Deception without debriefing; scraping personal data; a design whose
    control arm withholds something people need (\chref{ethics}) \\
\textbf{R}elevant
  & To practice, policy or subsequent research
  & Fails \reg{PhD Regs 2024}{\S14.1}: no connection to community need, the
    national research agenda, or an SDG \\
\bottomrule
\end{tabular}
\end{center}

\begin{pitfall}[Feasibility is checked with numbers, not optimism]
``Feasible'' is the criterion students assess by feel, and it is the one that
ends degrees. Do the arithmetic before committing: an a-priori power
calculation (\chref{sampling}) for a quantitative design, or an estimate of
interviews-to-saturation for a qualitative one, multiplied by the realistic
recruitment rate you can demonstrate --- not the one you hope for.

\medskip
If the arithmetic says the study is infeasible, that is a finding, arrived at
in week 10 rather than month 20. Portfolio piece P3 asks you to record exactly
such a decision, and Appendix~\ref{app:schedule} gives credit for it.
\end{pitfall}

\section{PICOC: Decomposing the Question}

PICOC breaks a question into the parts that drive both the design and, later,
the literature search~\cite{petticrew2006systematic}. It originates in clinical
research, and in computing the Comparison and Context elements do most of the
work.

\begin{center}
\small
\begin{tabular}{@{}L{2.4cm}L{4.4cm}L{7.1cm}@{}}
\toprule
 & \textbf{Element} & \textbf{Example (illustrative)} \\
\midrule
\textbf{P}opulation  & Who or what is studied
                     & Professional developers performing code review \\
\textbf{I}ntervention& The technique, tool or treatment
                     & An LLM-based review assistant \\
\textbf{C}omparison  & The alternative it is measured against
                     & Unassisted review; static analysis alone \\
\textbf{O}utcome     & The effect that matters
                     & Defects correctly identified; false positives raised \\
\textbf{C}ontext     & Setting, domain, constraints
                     & Industrial teams, safety-critical C code, time-boxed
                       reviews \\
\bottomrule
\end{tabular}
\end{center}

\begin{conceptbox}[The Comparison element is where most computing questions leak]
``Does an LLM assistant improve code review?'' has no C. Improve compared with
what --- no review at all, unassisted review, a linter, a second human
reviewer? Each comparison makes it a different study with a different
conclusion, and a paper that leaves C implicit will be attacked on exactly this
point.

\medskip
The same omission is what makes a literature search return thousands of
irrelevant hits: without C, there is no term to constrain the query
(\chref{slr}).
\end{conceptbox}

\begin{templatebox}[C1 --- Problem, gap and question]
Copy this into your portfolio (P1) and fill it in. One page, no more.

\medskip
\textbf{1. The problem, in the world.} Who has it, and what does it cost them?
\filllines{2}

\noindent\textbf{2. What is currently known.} Two or three sentences, with
citations.
\filllines{2}

\noindent\textbf{3. The gap.} What specifically is \emph{not} known? ``Little
research exists'' is not a gap; name the missing condition, population or
comparison.
\filllines{2}

\noindent\textbf{4. The question.} One sentence, ending in a question mark.
\filllines{1}

\noindent\textbf{5. Question type.} Exploratory / descriptive / relational /
causal / design / evaluative. \fillline{4cm}

\noindent\textbf{6. Why the answer matters.} What decision changes?
\filllines{2}
\end{templatebox}

\begin{templatebox}[C2 --- FINER and PICOC worksheet]
\textbf{PICOC}

\medskip
\begin{tabular}{@{}L{2.3cm}L{11.2cm}@{}}
Population   & \fillline{10.9cm} \\[7pt]
Intervention & \fillline{10.9cm} \\[7pt]
Comparison   & \fillline{10.9cm} \\[7pt]
Outcome      & \fillline{10.9cm} \\[7pt]
Context      & \fillline{10.9cm} \\[7pt]
\end{tabular}

\medskip
\textbf{FINER} --- for each, write the evidence, not a yes.

\medskip
\begin{tabular}{@{}L{2.3cm}L{11.2cm}@{}}
Feasible    & \fillline{10.9cm} \\[7pt]
Interesting & \fillline{10.9cm} \\[7pt]
Novel       & \fillline{10.9cm} \\[7pt]
Ethical     & \fillline{10.9cm} \\[7pt]
Relevant    & \fillline{10.9cm} \\[7pt]
\end{tabular}
\end{templatebox}

\section{From Construct to Measure}

A \term{construct} is the abstract thing you care about: code quality,
developer productivity, trust in an automated system, technical debt, usability.
None of these can be observed. What you observe is a \term{measure}, and the
distance between the two is where studies quietly go wrong.

\begin{definitionbox}[Operationalisation]
\term{Operationalisation} is the explicit statement of the rule by which an
abstract construct becomes a recorded observation.

\medskip
It is a claim, and it can be wrong. ``Developer productivity = commits per
week'' is an operationalisation; it is also, demonstrably, a bad one, because
it rewards splitting work and penalises the reading and design that good
development requires. \chref{measurement} treats validity of measures
properly; the point here is that the choice must be \emph{visible} so it can be
argued with.
\end{definitionbox}

\begin{worked}[Narrowing a topic in four rounds]
\textbf{Round 1 --- the topic.} ``Code review and software quality.''
Unbounded. Thousands of papers.

\medskip
\textbf{Round 2 --- add a population and a context.} ``Code review practices in
small software houses in Pakistan.'' Now bounded, but still not a question:
there is no outcome and nothing that could come out either way.

\medskip
\textbf{Round 3 --- add an outcome and a comparison.} ``Do teams in small
software houses that use checklist-guided code review find more defects than
teams using unguided review?'' A causal-comparative question, with a
falsifier: no difference, or fewer.

\medskip
\textbf{Round 4 --- confront feasibility, and revise honestly.} A between-teams
experiment needs many teams. Suppose you can access four. The power
calculation (\chref{sampling}) says four teams cannot detect any plausible
effect. Three legitimate moves:

\begin{tight}
  \item \textbf{Change the unit of analysis.} Randomise \emph{reviews} within
    teams rather than teams, using a within-subjects design
    (\chref{experiments}). Many more units, at the cost of a possible carryover
    effect that must then be controlled.
  \item \textbf{Change the question type.} Ask an exploratory question instead:
    how do reviewers in these firms actually decide what to comment on? Four
    teams is ample for that (\chref{qualitative}).
  \item \textbf{Change the outcome.} If defects found is too noisy at this
    scale, a construct closer to the intervention --- review coverage,
    comment specificity --- may be measurable with the units available.
\end{tight}

\medskip
All three are honest. Proceeding with four teams and reporting $p = 0.31$ as
``no significant difference'' is not: an underpowered null result is not
evidence of no effect (\chref{inference}).
\end{worked}

\subsection{Variables and their roles}

\begin{center}
\small
\begin{tabular}{@{}L{2.6cm}L{5.0cm}L{6.3cm}@{}}
\toprule
\textbf{Role} & \textbf{Definition} & \textbf{Example} \\
\midrule
Independent  & What you manipulate or contrast & Review style: checklist vs
                                                  unguided \\
Dependent    & What you measure as the outcome & Defects correctly identified \\
Moderator    & Changes the \emph{size} of the effect & Reviewer experience:
                                                       the checklist may help
                                                       novices more \\
Mediator     & Explains \emph{how} the effect happens & Attention to
                                                        error-prone constructs \\
Confounder   & Associated with both, and not controlled & Task difficulty, if
                                                          harder code happened
                                                          to get checklists \\
Control      & Held constant to remove it as an explanation & Same codebase,
                                                              same time box \\
\bottomrule
\end{tabular}
\end{center}

Confounders are the reason \chref{experiments} exists: randomisation is the
only technique that handles confounders you have not thought of.

\section{Hypotheses, and When Not to Have Them}

A \term{hypothesis} is a specific, testable prediction derived from the
question. Stating one commits you in advance, which is its whole value ---
see \chref{prereg} on why a prediction registered before the data is worth
more than the same prediction offered after.

\begin{center}
\small
\begin{tabular}{@{}L{3.2cm}L{11.2cm}@{}}
\toprule
$H_0$ (null) & Checklist and unguided review do not differ in mean defects
               identified. \\
$H_1$ (alternative, directional) & Checklist-guided review identifies more
               defects than unguided review. \\
\bottomrule
\end{tabular}
\end{center}

A directional hypothesis is stronger and should be used when theory genuinely
predicts a direction --- but it must be committed to beforehand, because
choosing the direction after seeing the data is a way of halving your p-value
without gaining any evidence (\chref{prereg}).

\begin{alertbox}[Exploratory studies should not invent hypotheses]
Not every study has hypotheses, and manufacturing them is a common way to
damage otherwise good work.

\medskip
An exploratory qualitative study is trying to discover what the relevant
categories \emph{are}; it cannot predict relations between categories it has
not yet found. A systematic mapping study characterises a literature. A design
science project asks whether an artefact can be built and what it teaches.
None of these needs a null hypothesis, and bolting one on invites an examiner
to ask why it was not tested properly.

\medskip
State research questions. Add hypotheses only where you are genuinely
predicting, in advance, a specific relation --- and then test exactly those,
reporting anything else you find as exploratory.
\end{alertbox}

\section{The Alignment Table}

This is the single most valuable page in a synopsis, and the one examiners turn
to first. It forces every question to have a plan, and exposes any question
that does not.

\begin{center}
\small
\begin{tabular}{@{}L{0.7cm}L{3.4cm}L{2.6cm}L{3.2cm}L{3.4cm}@{}}
\toprule
 & \textbf{Question} & \textbf{Data} & \textbf{Source} & \textbf{Analysis} \\
\midrule
RQ1 & How do reviewers decide what to comment on?
    & 12 interviews
    & 4 firms, purposive
    & Thematic analysis, two coders \\
RQ2 & Does a checklist change comment specificity?
    & 240 reviews
    & Within-subjects, randomised
    & Mixed-effects model, reviewer as random effect \\
RQ3 & Does the effect depend on experience?
    & As RQ2
    & As RQ2
    & Experience $\times$ condition interaction \\
\bottomrule
\end{tabular}
\end{center}

\begin{conceptbox}[How to use it as a diagnostic]
Read each row across and ask the four killing questions.
\begin{tightnum}
  \item Can this \emph{analysis} produce an answer to this \emph{question}? A
    thematic analysis cannot answer ``how much''; a regression cannot answer
    ``what forms does this take''.
  \item Does this \emph{data} exist, and can you get it? Write the number, and
    where it comes from.
  \item Does the analysis assume something the data cannot supply ---
    independence, when the same reviewer contributes 20 reviews?
  \item Is any row empty? An empty cell is not an oversight, it is an
    unanswerable question. Delete the question or fix the row.
\end{tightnum}
\medskip
An alignment table built after data collection is a formality. Built before, it
is the cheapest possible way to discover that your study cannot work.
\end{conceptbox}

\begin{pitfall}[Four questions that cannot be answered as written]
\textbf{The compound question.} ``How do developers use static analysis tools
and what are the barriers to adoption and how can these be overcome?'' Three
questions, and the third is not a research question at all --- it asks for a
recommendation, which is a conclusion, not a finding.

\medskip
\textbf{The question with no possible negative.} ``How can machine learning
improve intrusion detection?'' presupposes that it does. Ask whether, and under
what conditions.

\medskip
\textbf{The question already answered.} Detectable only by \chref{slr}, which
is why the review precedes the design.

\medskip
\textbf{The question whose construct cannot be measured in your setting.}
``Does technical debt reduce team morale?'' --- both constructs are contested,
neither has an agreed measure, and you have one semester. Either commit to
defending an operationalisation properly (\chref{measurement}), or choose
constructs with validated instruments.
\end{pitfall}

\begin{traditions}{the research question}
\tradEMP{The question names variables and a population, and its answer is a
quantity with an interval. It must be settled before data collection, because
after that it is no longer a prediction.}
\tradDSR{The question is a design problem: can an artefact meeting requirements
R be constructed for context C, and what generalisable design knowledge does
constructing it yield?}
\tradQUAL{The question is open and expected to evolve. What is going on here,
and how do participants understand it? Fixing it too early forecloses exactly
what the study exists to discover.}
\tradFORM{The question is a conjecture about a formal object: does this
property hold, what is the bound, is the problem decidable? It is settled by
proof, and no sample size applies.}
\end{traditions}

\begin{lab}[Build your own alignment table]
Using template C1 and C2, take your own topic through the four rounds of the
worked example. Then draw the alignment table with a row per research question.

\medskip
Now apply the backward pass from \figref{alignchain}: start at the analysis
column and read right to left. For at least one row you will find that the
analysis does not answer the question as stated. Fix the row --- by changing
the question, the data or the analysis --- and note in one sentence which you
changed and why. That sentence is the beginning of portfolio piece P3.
\end{lab}

\begin{checkpoint}
\begin{tightnum}
  \item Convert this into a PICOC question: ``I want to study whether
    containers help deployment.''
  \item A study asks ``Does pair programming improve code quality?'' and
    measures quality as lines of comments per function. State the
    operationalisation, and give the strongest objection to it.
  \item Why is a directional hypothesis chosen after seeing the data worth less
    than the same hypothesis chosen before, even though the sentence is
    identical?
  \item Give a research question in your own area for which stating a null
    hypothesis would be inappropriate, and explain why.
\end{tightnum}
\end{checkpoint}

\begin{chaptersummary}
\begin{tight}
  \item Topic, question, construct, variable, measure, analysis form a chain.
    Build it forwards; verify it backwards before collecting anything.
  \item Six question types, each admitting different designs. A causal question
    answered by a correlational study is a fatal mismatch, not a minor one.
  \item FINER rejects questions early and cheaply. Feasibility is settled with
    a power calculation or a saturation estimate, not with optimism.
  \item PICOC decomposes the question and later drives the search string. The
    Comparison element is the one computing questions most often omit.
  \item Operationalisation is a claim that can be wrong. Make it visible so it
    can be argued with.
  \item The alignment table --- question, data, source, analysis --- is the
    cheapest diagnostic in the course. An empty cell means an unanswerable
    question.
  \item State hypotheses only where you genuinely predict in advance.
    Exploratory work states questions instead.
\end{tight}
\end{chaptersummary}

\begin{reviewq}
  \item Easterbrook et al. argue method follows from question and stance. Give
    a case in computing where the same question could legitimately be answered
    within two different paradigms, and say what each answer would and would
    not establish.
  \item A reviewer objects that your operationalisation of ``maintainability''
    as cyclomatic complexity is invalid. Construct the strongest version of
    their objection, then the strongest defence available to you, and say what
    additional evidence would settle it.
  \item The alignment table exposes underpowered rows before data collection.
    Why, then, do so many published studies contain them? Give one structural
    and one psychological reason.
  \item Under what circumstances is it legitimate to change a research question
    \emph{after} data collection has begun? Distinguish carefully between the
    legitimate cases and HARKing (\chref{prereg}).
\end{reviewq}
